\documentclass[instructions.tex]{subfiles}
\begin{document}
 
\chapter{De la chasteté.}
 
\primo La chasteté est une vertu si belle, que ses ennemis ne peuvent s'empêcher de l'admirer. C'est un trésor plus précieux que tout l'or de l'univers\footnote{Omnis ponderatio auri non est digna continentis anima. Eccl. 26.}. Elle élève l'homme au-dessus de l'ange, dit saint Basile. Dans les anges, la pureté est sans combats; dans l'homme, elle est le fruit de sa fidélité. C'est pour cela que les anges ont tant de respect pour les personnes chastes, et que les démons les craignent.
 
Une âme livrée à l'impureté ne peut s'élever à Dieu, ni goûter les choses du ciel, parce qu'elle est comme ensevelie dans la chair. Une âme chaste au contraire s'unit facilement à Dieu, et Dieu se communique à elle si efficacement, qu'il n'y a point de graces plus spéciales que celles que Dieu donne aux âmes pures.
 
Jésus-Christ a honoré singulièrement cette vertu à sa naissance. Si un Dieu devait naître, dit saint Augustin, il ne pouvait naître que d'une vierge; si une vierge devait enfanter, elle ne pouvait enfanter qu'un Dieu. Il l'a honorée par sa conduite, se choisissant un précurseur vierge, une mère toujours vierge, un père nourricier vierge, un disciple favori vierge. Il n'a jamais permis qu'on donnât atteinte à la réputation de ses apôtres en cette matière, et quoiqu'il ait souffert parmi eux un avare et un incrédule, il n'en a souffert aucun qui n'ait été chaste.
 
\secundo Les délices que Dieu fait goûter aux âmes pures sont incompréhensibles. Non, dit saint Cyprien, il n'est point de plus grand plaisir que d'avoir vaincu la chair et la volupté.
 
Un impudique, au contraire, souffre dès cette vie un enfer anticipé, par les remords qui lui déchirent l'âme, et par la confusion qui l'accable. Il est vrai qu'il faut soutenir des combats pour conserver cette vertu. C'est une rose qu'on ne peut cueillir que parmi des épines. Mais que ces combats sont consolans, et qu'ils sont glorieux, puisque la chasteté a quelquefois plus de mérite et de gloire que le martyre! Les martyrs ont combattu peu de temps, mais les combats de la chasteté sont continuels. S'il faut du courage pour vaincre une fois les bourreaux, il en faut plus pour se vaincre toujours soi-même.
 
Il faut l'avouer, les âmes les plus chastes sont souvent celles qui ont le plus de combats à soutenir, qui sont les plus attaquées de tentations et de pensées affreuses. Faut-il s'en étonner? Quel mérite aurait-on si l'on était sans tentations? On ne peut être victorieux sans combats, ni couronné sans victoire. Dieu vous voit, âme fidèle; il voit le fond de votre cœur. Plus la tentation est violente, plus Dieu est près de vous et plus il récompensera votre courage. Vous ne pouvez résister sans la grace, parce que la continence est un don de Dieu. Demandez-la incessamment, en disant avec saint Paul: Infortuné que je suis! qui est-ce qui me délivrera de ce corps de mort? votre grace, ô mon Dieu!
 
\tertio On fait injure à la chasteté lorsqu'on la croit trop difficile; et l'on fait injure à la vérité lorsque, à l'exemple de l'impudique Luther, on la regarde comme impossible. Dire aux jeunes gens que la chasteté est une vertu qui ne convient ni à leur humeur, ni à leur âge, c'est les tromper et leur faire autant de tort qu'aux Israélites, quand on leur disait que la terre promise était un pays rempli de monstres; et ils n'y trouvèrent que des délices. Personne ne dit que la chasteté est impraticable, excepté les gens ignorans ou vicieux.
 
Pour être chaste, il ne faut être ni incivil ni affecté; mais il faut veiller sur son cœur, être réservé dans ses discours, modeste dans ses parures et dans ses manières. Si l'on doit être affable et complaisant, il ne faut pas que ce soit aux dépens de la pudeur. La complaisance devient un crime quand elle passe les bornes de la modestie. Par les ris et les cajoleries on souille son âme, et on se prépare des sujets de repentir\footnote{Risus dolore miscebitur. Prov. 14.}. Un air grave et sévère fait conserver à une fille ce que la complaisance lui fait perdre\footnote{Tamdiù virgo, quamdiù austera. SS. PP.}.
 
On n'est plus chaste quand on est coquette, qu'on aime celles qui le sont, et parce qu'elles le sont. On n'est pas même chaste, dit Tertullien, quand on a fait soupçonner de sa pudeur\footnote{Alterius suspicione violatur castitas. Tert.}. Si l'on aime une personne du sexe, que ce soit pour sa vertu et par un motif saint. Si vous ne l'aimez que pour sa beauté et ses agrémens, vous lui devenez dangereux, et votre cœur sera bientôt empoisonné par cet appât; comme le poisson, qui, en prenant un hameçon empoisonné, se tue lui-même, et empoisonne ceux qui en goûtent.
 
La chasteté est une vertu aussi fragile qu'elle est belle. C'est une brillante glace qu'un léger souffle ternit; un regard peut lui donner le coup de la mort. Ce coup est d'autant plus funeste, qu'on le reçoit sans y prendre garde. Veillez donc sur vous, et veillez en tout temps, dit Jésus-Christ.
 
\section{Résolutions}
 
\primo Puisque la chasteté plaît tant à Dieu, j'en aurai aussi moi-même une grande estime le reste de ma vie .

\secundo Je la demanderai souvent dans mes élévations de cœur.

\tertio J'éviterai avec soin tout ce qui pourrait me la faire perdre.

\quarto Et quand je m'apercevrai que je suis tenté, je détournerai d'abord mon esprit et mes yeux de l'objet qui me porte au mal; j'invoquerai le secours du ciel par quelque prière courte, mais fervente.

\prayer{Mon Dieu, vous voyez mes combats, ma foiblesse, et les dangers qui m'entourent; hâtez-vous de venir à mon aide.}
 
\end{document}
